% =============================================================================
% 03_neuro_oncology.tex -- Neuro-Oncological Challenges (~5-6 frames)
% =============================================================================

\section{Neuro-Oncological Challenges}

% ---- Frame 1 ----
\begin{frame}{Diffuse Gliomas: Biology and Mechanics}

  \begin{columns}[T]
    \column{0.55\textwidth}
    \begin{itemize}
      \item Infiltrative growth along white matter tracts
      \item Mass effect: compression of adjacent structures
      \item Peritumoral edema: altered tissue properties
      \item Blood-brain barrier disruption
      \item Grades II--IV (WHO 2021 classification)
    \end{itemize}

    \bigskip

    \textbf{Mechanical consequences:}
    \begin{itemize}
      \item Elevated intracranial pressure (ICP)
      \item Midline shift, ventricular compression
      \item Altered viscoelastic properties in tumor zone
    \end{itemize}

    \column{0.4\textwidth}
    % FIGURE SLOT: Freeview T1w with visible tumor (Google Slides)
    \begin{center}
      \fbox{\parbox{0.9\textwidth}{\centering\small
        \textit{GSlides:} T1w showing\\
        glioma mass effect\\[2pt]
        \texttt{FV-01\_t1w\_tumor.png}
      }}
    \end{center}
  \end{columns}

\end{frame}

% ---- Frame 2 ----
\begin{frame}{Surgical Context: Resection and Brain Shift}

  \begin{columns}[T]
    \column{0.55\textwidth}
    \textbf{During surgery:}
    \begin{itemize}
      \item Craniotomy $\to$ CSF drainage $\to$ brain shift
      \item Intraoperative deformation: 10--20\,mm
        \cite{Clatz2005}
      \item Tissue removal creates cavity
    \end{itemize}

    \bigskip

    \textbf{Post-operatively:}
    \begin{itemize}
      \item Resection cavity fills with CSF
      \item Surrounding tissue remodels
      \item Ventricles may expand
      \item Connectivity reorganizes
    \end{itemize}

    \column{0.4\textwidth}
    % FIGURE SLOT: pre-op vs post-op T1w comparison
    \begin{center}
      \fbox{\parbox{0.9\textwidth}{\centering\small
        \textit{GSlides:} Pre-op vs\\
        post-op T1w overlay\\[2pt]
        \texttt{FV-02\_preop\_postop.png}
      }}
    \end{center}
  \end{columns}

  \bigskip

  \alert{Key insight:} The pre-op $\to$ post-op deformation is
  \emph{empirically observable} from MRI --- a rare source of
  ground-truth displacement data for brain mechanics.

\end{frame}

% ---- Frame 3 ----
\begin{frame}{Why Standard Pipelines Fail}

  Pathological anatomy breaks assumptions:

  \bigskip

  \begin{tabular}{lll}
    \toprule
    \textbf{Pipeline step} & \textbf{Assumption} & \textbf{Glioma violation} \\
    \midrule
    FreeSurfer \texttt{recon-all} & Normal cortical folding &
      Tumor displaces cortex \\
    Atlas registration & Diffeomorphic to template &
      Missing tissue (cavity) \\
    Tractography & Coherent WM tracts &
      Infiltrated / disrupted fibers \\
    DTI fitting & Gaussian diffusion &
      Edema alters diffusion \\
    \bottomrule
  \end{tabular}

  \bigskip

  \textbf{Our approach:}
  \begin{itemize}
    \item \texttt{recon-all-clinical} (SynthSeg) for post-op sessions
    \item Automatic resection cavity masking for registration
    \item Cost-function masking in ANTs SyN
  \end{itemize}

\end{frame}

% ---- Frame 4 ----
\begin{frame}{Longitudinal Monitoring: What Changes?}

  Pre-operative $\to$ post-operative comparison reveals:

  \bigskip

  \begin{columns}[T]
    \column{0.48\textwidth}
    \textbf{Structural changes:}
    \begin{itemize}
      \item Resection cavity volume
      \item Ventricular expansion
      \item Cortical thickness alterations
      \item White matter tract disruption
    \end{itemize}

    \column{0.48\textwidth}
    \textbf{Quantifiable via MRI:}
    \begin{itemize}
      \item $\Delta$FA, $\Delta$MD maps (dMRI)
      \item Displacement field $\bfu(\bfx)$ (SyN)
      \item Jacobian $J(\bfx)$ (volume change)
      \item $\Delta$SC, $\Delta$FC matrices
    \end{itemize}
  \end{columns}

  \bigskip

  \alert{$\Rightarrow$} Each quantity has a direct mechanical
  interpretation: displacement, strain, material property change,
  connectivity rewiring.

\end{frame}

% ---- Frame 5 ----
\begin{frame}{Key References}

  \begin{itemize}
    \item \textbf{Aerts et al.\ (2018)} \cite{Aerts2018} --
      Brain dynamics in glioma patients using The Virtual Brain
    \item \textbf{Aerts et al.\ (2022)} \cite{Aerts2022} --
      BTC dataset: pre- and post-operative mpMRI, OpenNeuro ds002080
    \item \textbf{Zarodniuk et al.\ (2026)} \cite{Zarodniuk2026} --
      Mechanical compression induces neuronal apoptosis
    \item \textbf{Clatz et al.\ (2005)} \cite{Clatz2005} --
      Nonrigid registration for intraoperative brain shift
    \item \textbf{Miga et al.\ (2016)} \cite{Miga2016} --
      Model-updated image guidance for brain shift compensation
  \end{itemize}

\end{frame}
